\section{Sintaks Dasar: Statement, Komentar, dan Aturan Penulisan}

Program JavaScript terdiri dari urutan \textbf{statement}---instruksi yang dieksekusi satu per satu dari atas ke bawah. Setiap statement biasanya diakhiri dengan titik koma (\code{;}), meskipun JavaScript memiliki fitur \textit{automatic semicolon insertion} (ASI) yang menambahkan titik koma secara otomatis di beberapa kasus. Praktik terbaik: \textbf{selalu tulis titik koma secara eksplisit} untuk menghindari bug yang sulit dilacak \cite{jsInfo}.

JavaScript bersifat \textbf{case-sensitive}: \code{myVariable} dan \code{myvariable} adalah dua identifier berbeda. Nama variabel dan fungsi harus dimulai dengan huruf, underscore (\code{\_}), atau dollar sign (\code{\$}), dan tidak boleh menggunakan kata kunci reserved (\code{if}, \code{for}, \code{function}, dll.). Konvensi penamaan yang umum adalah \textbf{camelCase} untuk variabel dan fungsi (\code{firstName}), \textbf{PascalCase} untuk class (\code{UserProfile}), dan \textbf{UPPER\_SNAKE\_CASE} untuk konstanta (\code{MAX\_SIZE}) \cite{mdnJsGuide}.

\begin{webcode}[language=JavaScript, caption={Statement, komentar, dan penamaan}]
// Komentar satu baris
/* Komentar
   multi-baris */

let firstName = "Budi";      // camelCase
const MAX_RETRIES = 3;        // UPPER_SNAKE_CASE
let isLoggedIn = true;        // boolean

console.log(firstName);       // Output: Budi
console.log(typeof isLoggedIn); // Output: boolean
\end{webcode}

\begin{konsep}
\textbf{Strict Mode.} Menambahkan \code{"use strict";} di awal file atau fungsi mengaktifkan mode ketat yang melarang beberapa praktik rawan error: variabel tanpa deklarasi menjadi error, assignment ke properti read-only gagal, dan beberapa sintaks yang deprecated dilarang. ES6 modules secara otomatis berjalan di strict mode. Gunakan strict mode untuk kode yang lebih aman dan mudah di-debug \cite{jsInfo}.
\end{konsep}

